% Part: first-order-logic
% Chapter: sequent-calculus
% Section: soundness-identity

\documentclass[../../../include/open-logic-section]{subfiles}

\begin{document}

\olfileid{fol}{seq}{sid}

\olsection{\usetoken{S}{identity}-সহ বিশুদ্ধতা}

\begin{prop}
অভিন্নতার প্রারম্ভিক সিকোয়েন্ট ও বিধিসহ $\Log{LK}$ বিশুদ্ধ।
\end{prop}

\begin{proof}
${} \Sequent \eq[t][t]$ আকারের প্রারম্ভিক সিকোয়েন্টগুলি বৈধ, কারণ
প্রত্যেক !!{structure}~$\Struct M$-এর জন্য $\Sat{M}{\eq[t][t]}$। (লক্ষ
করো, পদ~$t$-কে আমরা বদ্ধ ধরে নিচ্ছি, অর্থাৎ এতে কোনো চলরাশি নেই; তাই
চলরাশি-আরোপ প্রাসঙ্গিক নয়।)

ধরা যাক, কোনো !!a{derivation}-এর শেষ অনুমানটি $=$। তাহলে পূর্বধারণা
$\eq[t_1][t_2], \Gamma \Sequent \Delta, !A(t_1)$ এবং সিদ্ধান্ত
$\eq[t_1][t_2], \Gamma \Sequent \Delta, !A(t_2)$। কোনো
!!a{structure}~$\Struct M$ বিবেচনা করি। দেখাতে হবে যে সিদ্ধান্তটি
বৈধ, অর্থাৎ $\Sat{M}{\eq[t_1][t_2]}$ ও $\Sat{M}{\Gamma}$ হলে হয়
কোনো $!C \in \Delta$-এর জন্য $\Sat{M}{!C}$, নয়তো $\Sat{M}{!A(t_2)}$।

আরোহের অনুমান অনুসারে পূর্বধারণাটি বৈধ। এর অর্থ,
$\Sat{M}{\eq[t_1][t_2]}$ ও $\Sat{M}{\Gamma}$ হলে হয় (a) কোনো $!C
\in \Delta$-এর জন্য $\Sat{M}{!C}$, অথবা (b) $\Sat{M}{!A(t_1)}$।
(a) ক্ষেত্রে কাজ শেষ। (b) ক্ষেত্রটি বিবেচনা করি। এমন একটি চলরাশি-আরোপ
$s$ নিই যেন $s(x) = \Value{t_1}{M}$।
\olref[syn][ass]{prop:sentence-sat-true} অনুসারে
$\Sat{M}{!A(t_1)}[s]$। যেহেতু $\varAssign{s}{s}{x}$, তাই
\olref[syn][ext]{prop:ext-formulas} অনুসারে $\Sat{M}{!A(x)}[s]$। আবার
$\Sat{M}{\eq[t_1][t_2]}$ বলে $\Value{t_1}{M} = \Value{t_2}{M}$, এবং
তাই $s(x) = \Value{t_2}{M}$। পুনরায়
\olref[syn][ext]{prop:ext-formulas} প্রয়োগ করে
$\Sat{M}{!A(t_2)}[s]$-ও পাই। শেষে
\olref[syn][ass]{prop:sentence-sat-true} অনুসারে $\Sat{M}{!A(t_2)}$।
অভিন্নতার দ্বিতীয় বিধির ক্ষেত্রটি পদদুটির ভূমিকা অদলবদল করে একইভাবে
প্রমাণিত হয়।
% BN-SRC-032 TARGET-CORRECTION-BEGIN
\par\smallskip\noindent\textnormal{\small\textbf{উৎস-সংশোধন (BN-SRC-032):} আগের অংশে অভিন্নতার দুটি অনুমানবিধি দেওয়া হলেও হিমায়িত ইংরেজি বিশুদ্ধতা-প্রমাণটি শুধু প্রথম অভিমুখটি যাচাই করে শেষ হয়েছে। বাংলা প্রমাণে দ্বিতীয় অভিমুখটির সমমিত যুক্তি স্পষ্ট করে একটি বাক্য যোগ করা হয়েছে; কোনো বিধি বা সূত্র বদলানো হয়নি।} % BN-SRC-032 TARGET-CORRECTION-END
\end{proof}

\end{document}
